\section{Lacunas Clínicas: A Barreira Humana na Translação Digital}

Mesmo que a engenharia de software resolva todas as imperfeições da renderização e simulação, a adoção em larga escala de gêmeos digitais colidirá irremediavelmente com fatores humanos, organizacionais e culturais profundamente arraigados na prática clínica. O ambiente odontológico é historicamente pautado no empirismo, na intuição visual e na destreza manual; a transição para um modelo analítico-computacional exige mais do que novos softwares, exige um letramento digital (\textit{digital literacy}) substancial.

\subsection{Ceticismo e a Síndrome da ``Caixa-Preta''}

A resistência à adoção de tecnologias não é um fenômeno irracional, mas frequentemente uma resposta prudente à incerteza. Profissionais altamente experientes, que desenvolveram uma intuição clínica afiada ao longo de décadas, veem com justificável ceticismo os modelos computacionais preditivos. Esses algoritmos, muitas vezes fundamentados em redes neurais de aprendizado profundo (Deep Learning), operam como ``caixas-pretas'': fornecem um plano de tratamento ótimo sem explicitar claramente o raciocínio subjacente. Este déficit de explicabilidade (\textit{Explainable AI - XAI}) fere o princípio da autonomia clínica, criando uma barreira de confiança \cite{mdpi2025digital}. Sem entender "o porquê" da recomendação, o profissional se recusa a assumir a responsabilidade civil pelo resultado.

\subsection{A Fricção no Fluxo de Trabalho (Workflow Integration)}

Um erro comum das empresas de tecnologia (\textit{HealthTechs}) é conceber o gêmeo digital como um ``produto adicional'' no portfólio da clínica, ignorando a grave fricção que ele impõe à rotina. O tempo despendido para calibrar malhas, corrigir segmentações automáticas imperfeitas e ajustar condições de contorno frequentemente ultrapassa o tempo da consulta presencial. Se o ganho preditivo não for superior à perda de produtividade, a tecnologia é rapidamente abandonada. Redesenhar o fluxo de trabalho exige uma engenharia de processos (\textit{Business Process Reengineering}) para a qual a imensa maioria dos consultórios independentes não está capacitada.

\subsection{O Custo Cognitivo do Aprendizado}

A curva de aprendizado para a manipulação de softwares odontológicos avançados é íngreme. O cirurgião-dentista não foi formado, em suas bases acadêmicas, para compreender conceitos de malhas volumétricas, anisotropia de materiais ou tensores de estresse de Von Mises. A ausência de interfaces de usuário (UI) intuitivas e centradas no profissional obriga o clínico a atuar como um engenheiro assistente. A literatura de difusão de inovação aponta que, até o presente, a tecnologia atende primordialmente a uma vanguarda de \textit{early adopters} (em torno de 15\% do mercado). Para cruzar o abismo (\textit{crossing the chasm}) rumo à maioria pragmática, os sistemas precisarão evoluir para um paradigma de automação invisível.
